\scanpage{33}
En fait, on trouve un \unsure{diagramme|diagr.me|diagr.}
ess. commutatif
\diagram{\begin{tikzcd}[ampersand replacement=\&,column sep=0.45em,row sep=3em]
\&[-1em]\mathcal A \arrow[dl,"{\underline{\mathrm{Id}}}"'] \arrow[drrr,no head,"{\underline{\mathrm{id}}}"'] \& \subset \& C \arrow[dll,"{\simeq}"',"{v}"] \arrow[drr,no head,"{v'}","{\simeq}"'] \& \& \& \\
\mathcal A \& {\subset\widehat{\mathcal A}_{\mathcal T}} \& {\subset\widehat{\mathcal A}} \& \& \mathcal A \arrow[r,hook] \& {(\widehat{\mathcal A^\circ}_{\mathcal T'})^\circ} \arrow[r,hook] \& {(\widehat{\mathcal A^\circ})^\circ}
\end{tikzcd}}
(resp. en $\widehat{\mathcal A}_{\mathrm{ab}}$ et $\widehat{\mathcal A^\circ}_{\mathrm{ab}}$)

On trouve en particulier une
équivalence remarquable
\[
 \widehat{\mathcal A}_{\mathcal T}\simeq
 (\widehat{\mathcal A^\circ}_{\mathcal T'})^\circ
\]
qui peut se définir intrinsèquement \inserted{en termes de $\mathcal A,\mathcal T,\mathcal T'$,} chaque fois qu'on a une
sous-cat. $D$ de $\widehat{\mathcal A}$ (resp. une sous cat.
additive de $\widehat{\mathcal A}_{\mathrm{ab}}$), stable \inserted{dans $\widehat{\mathcal A}$} par
$\varprojlim$ de type $\mathcal T$, et \inserted{où existent les} $\varinjlim$ de type $\mathcal T'$,
(pas néc. les $=\mathrm{s}$ que dans $\widehat{\mathcal A}$ !)
et dont chaque objet est à la fois
$\varprojlim$ \inserted{de type $\mathcal T$} dans $C$ de foncteurs repr.,
et $\varinjlim$ \inserted{de type $\mathcal T'$} dans $C$ de tels foncteurs.

Si $\mathcal A$ est réduit à un seul objet, \inserted{$A$} et
dans le cas \unsure{resp$^\circ$|resp.} ($C$ additive)
$\mathcal T=\mathcal T'=$ ens. catégories discrètes finies,
on trouve une antiéquivalence
entre modules libres de t. f à g. et
% archive sheet number: 32
